\chapter{Paper versus code: every difference}\label{ch:differences}

This chapter is the one to bring to a meeting. It compares strict CSC-LISTA against Eq.~20 and Appendix H
of Liu et al. It is more complete than the table in \code{README_CSC_LISTA.md}.

\section{What matches the paper}

\begin{center}
\renewcommand{\arraystretch}{1.3}
\begin{tabularx}{\textwidth}{lX}
\toprule
Item & Status \\
\midrule
Layer update & Eq.~20, term by term. See Chapter~\ref{ch:csc_lista} \\
Filter per layer & $w^{(k)}$ learned separately per layer \\
Threshold & One $\theta^{(k)}$ per layer, shared by all atoms \\
Output & Codes convolved with the original dictionary, as in Appendix H \\
Loss & MSE between output and clean image, as in Appendix H \\
Order of work & Dictionary first, then the network with the dictionary fixed \\
\bottomrule
\end{tabularx}
\end{center}

\section{What is necessary for our problem}

\begin{center}
\renewcommand{\arraystretch}{1.3}
\begin{tabularx}{\textwidth}{>{\raggedright\arraybackslash}p{3.1cm}>{\raggedright\arraybackslash}X>{\raggedright\arraybackslash}X}
\toprule
Item & Paper & Ours, and why \\
\midrule
Task & Denoising, $b=t+$noise & Deblurring, $b=h*t+$noise. This is our problem \\
Dictionary in the data term & Small learned $d_m$, used directly & $d_m=h*g_m$, long blurred filters. This is how the blur gets into Eq.~20. No published result covers this use \\
Normalisation & $\lVert d_m\rVert_2=1$ & $\lVert g_m\rVert_2=1$, so the blurred $d_m$ are not unit length \\
Data & 400 natural images & Simulated fields, 512 pixel crops \\
Boundary & Non-circular, $x$ larger than $b$ (Eq.~18) & FFT with zero padding, codes on the padded canvas, and a window mask on the residual. Same purpose, different method \\
Layers & 5 & 16, 32 and 64, because the blur is far harder than the denoising task \\
\bottomrule
\end{tabularx}
\end{center}

\section{What is our own engineering choice}

\begin{center}
\renewcommand{\arraystretch}{1.3}
\begin{tabularx}{\textwidth}{>{\raggedright\arraybackslash}p{3.1cm}>{\raggedright\arraybackslash}X>{\raggedright\arraybackslash}X}
\toprule
Item & Paper & Ours \\
\midrule
Size of $w$ & Same size as $d_m$ & 155 pixels, shorter than the long blurred $d_m$. Trials showed little gain from larger \\
Dictionary learning & Algorithm 1 of Liu et al.\ 2018 & Alternating FISTA and Adam \\
Atoms & $M=64$, $D=7$ in its experiments & 32 atoms of 75 pixels \\
Start of $w$ & Not specified for Eq.~20 & ISTA values, $d_m/L$, truncated \\
Start of $\theta$ & Not specified for Eq.~20 & $\lambda/L$, with $\lambda=0.3$ \\
Optimiser & Not specified in Appendix H & Adam, with a rate 3\% of each parameter's scale \\
Rate schedule & Not specified & Cosine decay to zero \\
Gradient clipping & Not mentioned & Norm limited to 1.0 \\
Threshold sign & $\theta$ is a threshold, so $\ge0$ & \code{clamp_min(0)} in the layer, and a clamp at 0 after each step \\
Loss region & Not stated & Middle of the crop, 64 pixel margin dropped \\
Checkpoint & Not stated & Best simulated validation score \\
\bottomrule
\end{tabularx}
\end{center}

\section{Additions that strict runs removed}

\begin{center}
\renewcommand{\arraystretch}{1.3}
\begin{tabularx}{\textwidth}{lXl}
\toprule
Item & What it did & Used in \\
\midrule
Threshold floor & Kept $\theta$ at or above $\lambda/L$ & Sparse runs only \\
Relative loss & Replaced MSE with per-crop relative $L_2$ & Relative runs only \\
\bottomrule
\end{tabularx}
\end{center}

\begin{defend}[The short version to memorise]
``The layer, output and loss are from Eq.~20 and Appendix H of Liu et al. Three things are necessarily
ours. We deblur rather than denoise, so the dictionary is $h*g$. We use FFT convolution with padding
and a mask. And our dictionary learning is a substitute for theirs. Everything else is standard training
practice, listed in full.''
\end{defend}

\begin{careful}
A note on the paper's own text. Section 5 says the dictionary came from Algorithm 1 of Liu et al.\ 2018, while
Appendix H says it was obtained ``by solving (24)'', which is the analytic weight problem. The second looks
like loose wording, but the two statements do not match. If you are asked, quote Section 5.
\end{careful}

\section{Which code actually ran}

Each server run used a frozen copy of the scripts, so later edits could not change it. The local \code{scripts/} folder,
which this booklet quotes, is byte-identical to the frozen copy for FISTA-Net run 2, and FISTA-Net run 1 used the same code before the option for start
values was added. The Job 3 and earlier CSC-LISTA runs used slightly older copies. They differ only in code those runs did not use, such as the momentum and Lasso
options, so the Eq.~20 layer quoted here is the one that ran. The FISTA depth scores also came from a frozen copy,
identical to \code{scripts/score_fista_depths.py}.

\section{How the code was checked}

Fifteen unit tests in \code{tests/test_csc_lista.py} pass on the server. These are the important ones:

\begin{center}
\begin{tabularx}{\textwidth}{lX}
\toprule
Test & What it proves \\
\midrule
\code{test_forward_and_adjoint_are_transposes} & The prime in Eq.~20 is correct \\
\code{test_bound_is_the_operator_norm} & $L$ is the true largest eigenvalue \\
\code{test_full_width_start_matches_ista} & The untrained network is ISTA \\
\code{test_residual_ignores_the_padding} & The mask works \\
\code{test_centred_filter_keeps_position} & Filters do not shift the image \\
\code{test_training_lowers_loss_and_reaches_every_layer} & Gradients reach every layer \\
\bottomrule
\end{tabularx}
\end{center}
